\section{The Propers: Detailed Commentary}
\zlabel{triptych:concise:commentary:start}

\subsection{What changes when the father goes home?}
The father asks Jesus to make a journey; Jesus gives him a journey to
make. This reversal exposes the difference between trusting a healer
and understanding whose power heals. Gregory's question is exact: why
would someone ask Christ to save his son unless he already believed
Christ capable of doing so? The request shows faith. Its limitation
lies in the requested means. The father thinks the Lord must be bodily
present beside the child, whereas divine power is nowhere absent.
Christ grants the desired life while correcting the restriction placed
upon his action. Mercy and instruction occur together.
\footnote{Gregory, \work{Homilies on the Gospels}, II.28.1.}

The father's repeated plea makes this correction more than an abstract
lesson about omnipresence. He fears death, and urgency binds his hope
to the means he knows. Christ does not require him first to compose
a sufficient account of divine power. The granted life itself reveals
the inadequacy of his first estimate. The father is not asked to care
less about the child; he is asked to accept a greater healer. His
departure is therefore neither resignation to death nor abandonment
of his responsibility. The way he cares has changed because of the
word he has received.

Augustine reads the rebuke through the earlier Samaritan encounter.
They heard Christ and believed; the Galileans' reception has followed
signs already seen. He judges the ruler's approach as cold or not yet
believing and takes the eventual statement of belief seriously as a
change. Gregory's explanation makes the father's original faith
defective; Augustine's diagnosis is harsher. Neither can simply be
made to say the other's words. Their agreement concerns Christ: he
knows the petitioner and is not confined to the visible conditions the
petitioner sets.\footnote{Augustine, \work{Tractates on John}, XVI.1--5.}

The two declarations of belief in John~4:50 and 53 show why this dispute
is worth hearing. At the first, trust is sufficient to govern departure.
At the second, the recovered child and matching hour are known, and the
household is included. The narrative distinguishes these moments without
requiring every person's faith to pass through a universal sequence of
identical stages. Augustine's comparison with Thomas commends belief
without demanding sight. It does not turn the servants' report into
an embarrassment. Testimony received after obedience and a visible
condition imposed before obedience play different parts in the story.

This distinction joins the personal and communal readings. The father's
departure answers the question of what trust now makes possible. The
household's belief answers the question of whom that mercy can gather.
He sought help on behalf of another, and those serving his household
bring him knowledge he does not possess. Gregory's comparison with the
centurion's servant adds a pointed moral consequence: Christ's regard
does not follow the social hierarchy that prizes a ruler above a
servant. The centurion belongs to Gregory's comparison, not to this
healing narrative.\footnote{Gregory, II.28.2--3.}

Ephesians' mutual submission puts that consequence into daily practice.
Chrysostom compares compulsory service for one master with friends
freely serving one another. Even a master owes provision and service
to those beneath him. The inward song cannot be detached from this
outward obligation without dropping the Epistle's final verse. The
Gospel's growth in faith and the Epistle's reciprocal conduct both
unsettle the supposition that receiving honor exempts a person from
attending to others.\footnote{Chrysostom, \work{Homilies on Ephesians},
XIX, on 5:19--21 and the concluding exhortation.}

\subsection{Does the need for mercy make the sufferer guilty?}
The Introit's confession and the Gospel's illness can be heard in one
Mass without becoming a causal explanation of each other. Azarias and
his companions enter the furnace because they refuse idolatry. Yet
Azarias confesses sin in the plural. Jerome explains that the faithful
youths speak in the person of their people. Their solidarity enlarges
the voice of confession; it does not reverse the narrative's account
of why they suffer. Nothing in John accuses the child of bringing his
illness upon himself.\footnote{Jerome, \work{Commentary on Daniel}, I,
at 3:23--29, especially 3:29, PL~25, cols.~639--640.}

This matters differently to the petitioner and to the assembly.
For the petitioner, confession admits need without requiring a
self-sufficient claim to deserve help. For the assembly, it prevents
the faithful from treating their own upright acts as grounds for
separation from a sinful people. Both belong before God as receivers
of mercy. Neither is authorized to pronounce an unrevealed judgment
on another person's wound. The youths' refusal of idolatry also
survives the confession: communal speech does not make the king's
violence good or erase particular responsibility.

The furnace's earlier scene gives this distinction another depth.
Before entering the furnace, the youths confess that God can deliver,
but refuse idolatry even if deliverance does not come (Dan.~3:17--18).
Their fidelity rests on who God is, not on a bargain for rescue.
John's father receives an explicit assurance that his son lives;
the youths declare fidelity without making the same outcome a
condition. Together these literal contexts show why trust cannot
be reduced to expecting one's desired result. They also keep the
assembly's corporate penitence from becoming a demand that each
afflicted person demonstrate faith by confidently predicting a cure.

The Collect asks pardon and peace so that service may become possible.
Schuster describes the heart distracted by passions and remorse;
forgiveness restores peaceful attention. That is a stronger gift than
temporary relief from an uncomfortable feeling. A mind released for
service can take the next faithful step, as the father does, or join
others in work that resentment has made impossible. The distinction
between the two readings is therefore an accent, not a conflict:
trust concerns action freed from the need to control the outcome;
thanksgiving concerns common life freed from the need to make one's
own advantage its measure.\footnote{Schuster, III, p.~175.}

Chrysostom's account of redeeming time clarifies the cost. Evil days
do not mean that time, God's creation, has become evil in its substance.
They mean that believers face conduct and circumstances demanding
discernment. Prudence may require relinquishing a lesser advantage
instead of answering hostility in kind. It is neither inactivity nor
the deliberate search for provocation. Fidelity retains an object
more important than winning the immediate encounter. Thus the Collect's
peace and the Epistle's wisdom give trust a durable moral shape beyond
the exceptional crisis of a sick child.
\footnote{Chrysostom, XIX, exposition of 5:15--17.}

\subsection{Whom does the open hand require to act?}
The Gradual praises provision for every living creature. The Gospel
tells of one restored child. Read together, they provoke a difficult
question: what of those still hungry, still ill, still waiting? The
words ``in due season'' lead Augustine to the physician's knowledge.
A patient may rightly desire a good without knowing when or how it
should be given. Delay need not establish abandonment. Augustine
encourages the person who asks to keep trusting the giver.
\footnote{Augustine, \work{Expositions of the Psalms}, Ps.~144,
\S\S14--17; NPNF heading Ps.~145.}

His answer does not stop at a promise that every earthly request will
eventually succeed on its original terms. He directs desire from
health, family and other benefits toward God himself. The physician's
greatest gift exceeds a recovered advantage within mortal life.
Resurrection and the final enjoyment of God let hope survive a temporal
answer different from the one desired. The child's cure remains a real
good; it is not the guarantee that every petitioner will receive the
same bodily outcome. The Offertory's tears and the Communion's
humiliation give the still-afflicted a voice within this very Mass.

Bellarmine asks the corresponding question of the person who already
has resources. God's universal generosity exposes human hoarding and
misuse. To praise an open divine hand while refusing to open one's own
is to set conduct against confession. His account presses rulers and
possessors toward generosity. Augustine's physician encourages the
person awaiting help; Bellarmine's open hand charges the person able
to help. These are complementary addresses. A comfortable listener
cannot use patience as the answer to a duty that belongs to himself.
Bellarmine's further conjectures about fault or improvidence supply
no diagnosis of any particular hungry person.
\footnote{Bellarmine, \work{Commentary on the Psalms}, Ps.~144:14--17.}

Schuster unfolds yet another relation: nourishment in this Mass leads
from gifts received under faith to Eucharistic food and finally to
vision. This sacramental reading does not replace the psalm's literal
praise of the Creator who feeds living things. It joins bodily need,
the gift of Christ and the desire for God without making them the
same object. The personal reading stresses the growth of the
petitioner's desire; the communal reading stresses the responsibility
of those sharing God's gifts. The anagogical end of both is God,
and that end strengthens rather than dismisses care for the neighbor.
\footnote{Schuster, III, pp.~175--176.}

The difference among these witnesses is consequential for the four
senses. Literally the chant concerns living creatures and their food;
Bellarmine's moral appeal asks what possessors do with those gifts.
Schuster's sacramental reading looks to Christ feeding his people,
and his movement toward vision gives that feeding an anagogical end.
Augustine's desire for God prevents the end from becoming merely an
unending supply of the benefits presently sought. These relations
allow hunger to remain a practical claim on charity while hope seeks
more than what charity can distribute. The question of delay is not
answered simply by turning every hungry creature into a symbol.

\subsection{How can thanksgiving include the captives' tears?}
Chrysostom refuses to confine thanks to prosperity. His exposition of
Ephesians includes poverty, sickness and anguish: faith must answer
God's goodness when its benefit is not yet understood. The nineteenth-
century NPNF editorial note~404 explicitly restricts ``all things'' to
blessings, against this wider reading. The note is not another sentence
of Chrysostom's homily. His demand remains more searching because
present satisfaction cannot be its condition.
\footnote{Chrysostom, XIX, exposition of 5:20; NPNF I.13, editorial
note~404, reporting Meyer's narrower judgment.}

Yet Chrysostom locates evil in human conduct rather than created days.
The distinction prevents thanks amid affliction from becoming approval
of the act that injures a neighbor. The Introit confesses wrongdoing
and the Secret seeks its removal. A praise that celebrated vice would
contradict these prayers. The Gospel's father seeks relief with urgency;
the literal value of the child's life gives bodily suffering its full
weight beside Chrysostom's rhetoric about the greater evil of sin.
The assembly may trust, seek help and accompany grief together without
claiming to know why every wound has been allowed.

Augustine approaches the same difficulty through what the exiles
remember. Babylon's currents signify passing loves whose movement
cannot provide an abiding home. Sitting beside them becomes a figure
of humility and resistance: the pilgrim recognizes captivity without
consenting to be swept away. But tears alone do not determine which
home he desires. One can grieve a worldly loss and still want nothing
beyond its replacement. To weep remembering Zion is to measure the
present by a good Babylon cannot supply. The object of grief changes
the direction of hope.\footnote{Augustine, Ps.~136, \S\S2--4
(NPNF heading Ps.~137).}

The full psalm's captors demand entertainment from the defeated.
Refusing that performance does not mean forgetting sacred song.
Augustine ends his exposition by calling the pilgrims to sing to one
another. Their music serves a shared journey toward God rather than
the amusement of those who mock their home. Ephesians similarly joins
address to one another with song to the Lord. Chrysostom asks for
attention within the singer; Augustine asks which love and destination
the singing sustains. Their readings explain how a community can
praise without demanding that grief first disappear.
\footnote{Augustine, Ps.~136, \S\S5--9,13; Chrysostom, XIX, on 5:19.}

Bellarmine's distinction between outward responsibilities and inward
captivity keeps this pilgrimage close to ordinary duty. One may
exercise authority or manage a household while regarding the task
as service rather than a possession to admire. Longing for Zion
does not require leaving dependents uncared for. The Alleluia's
prepared heart belongs to this freedom. Schuster reads its liturgical
address, \emph{tibi, gloria mea}, as God being the soul's glory; the
biblical wording rendered ``with my glory'' is different. Praise can
then remain steady without measuring itself by human approval.
\footnote{Bellarmine, Ps.~136:1--4; Schuster, III, p.~176. Augustine's
brief Ps.~107.1--2 (NPNF Ps.~108) identifies reused psalm passages;
it does not give Schuster's explanation of this liturgical address.}

The biblical Psalm~107 also contains threatened dependence and the
insufficiency of human help. Its ready heart does not belong to a
community immune from danger. That context joins Gregory's correction
of the ruler's estimate of power to Augustine's account of the loves
that carry a pilgrim away. Readiness is not possession of sufficient
earthly means. It is a direction of the heart toward the one whom
those means cannot replace. The liturgical address names praise's
object; the wider psalm explains why such praise can coexist with
need. Schuster's emphasis and the biblical context thus meet without
making his explanation of the chant Augustine's gloss on the verse.

The literal Babylonian captivity and its loss of Jerusalem remain
distinct from Augustine's allegorical account of the pilgrim Church.
His spiritual reading of the full psalm's violent close directs its
force against evil desires in their beginnings, not against persons.
That final verse is outside the Offertory. The appointed lament
gives this assembly tears; Christian hope gives those tears an end
beyond the recovery of every present possession.

\subsection{What does the medicine heal, and where does it lead?}
The Secret asks these mysteries to give heavenly medicine and expel
the vices of the heart. John describes a child's bodily recovery;
the prayer names the healing of disordered love. The ailments are
not interchangeable. Schuster's Eucharistic exposition locates the
remedy in Christ's gift, which the faithful receive because they
need healing. The same pride that resists obedience also refuses
mutual service; the same avarice that corrupts desire withholds food.
The personal and communal consequences therefore meet at their root,
without reducing the sacrament to a means of social convenience.
\footnote{Schuster, III, p.~177.}

The Communion's request to remember the word expresses desire for
the fulfilment of a promise. Augustine insists that God has not
forgotten. The promise itself gave the servant hope; the servant's
prayer now asks its accomplishment. Humiliation may include penitence,
trial or persecution. Comfort is possible within that lowliness,
and Augustine carries the Church's hope toward resurrection. The
biblical verse's concluding assertion that the word gives life is
context, omitted from the chant. Schuster's identification of the
Word with Christ supplies the further liturgical and Christological
accent: the Church's hope rests personally in him.
\footnote{Augustine, Ps.~118, \S\S50--56 (NPNF heading Ps.~119);
Schuster, III, p.~177.}

In Augustine's following exposition, God's statutes become songs in
the house of pilgrimage, and remembrance of the Lord endures in the
night. These images come from the psalm's context after the appointed
Communion, not from extra words in the chant. They show why promise
and remembered home belong together. The servant who asks fulfilment
does not have to abandon obedience while waiting; the pilgrims who
long for Zion do not have to abandon song. Remembrance sustains
conduct and shared worship across a period in which their final
object has not yet been possessed.\footnote{Augustine, Ps.~118,
\S\S53--56, on the context including verse~54.}

The two readings therefore locate comfort with different emphases.
The healing-word reading follows the promise that lets a petitioner
act before confirmation. The exile reading follows the promise that
lets a people keep their allegiance through continuing trial. Neither
requires comfort to mean that pain has ended. Augustine's Church can
be comforted under persecution, and Schuster's sacramental food can
sustain faith before vision. Their final hope is not a more persuasive
feeling about present circumstances but the fulfilment God gives.

The Postcommunion's request for worthiness might seem to reverse
the medicine's generosity. Its grammar resolves the difficulty:
God is asked to make the faithful obedient. They do not declare a
store of independent righteousness with which to buy his gifts.
Augustine's account of grace in Psalm~118 and Schuster's account of
habitual docility make the same distinction between receiving help
and claiming to have supplied its price. Obedience matters because
reception is to bear fruit. For the father-like petitioner this is
the next faithful step; for the pilgrim community it is the conduct
that endures after the singing ends.
\footnote{Augustine, Ps.~118, \S56; Schuster, III, pp.~177--178.}

In the commemorative branch, the Maternity Collect confesses that
the eternal Word really took flesh from Mary. Augustine's conclusion
to John XVI contemplates the Creator made man within the people he
made. Christ's power at a distance does not render his humanity
unreal, nor does his human birth diminish his divine power. The
Marian Secret asks present and perpetual peace through mercy and
intercession. Its Postcommunion asks cleansing and fellowship in
heavenly remedy; Lasance--Walsh's historical English identifies that
remedy personally with Christ. That explicit identification is the
translators' rendering, not an additional title in the Latin.
\footnote{Augustine, \work{Tractates on John}, XVI.7; Maternity prayers,
\work{Missale Romanum} (1962), pp.~685--686; Lasance--Walsh,
\work{The New Roman Missal}, pp.~1233,1235--1236.}

The Trinity Preface gives both readings their final direction.
The assembly's unity is not itself the object of worship. Thanks
filled with the Spirit, offered to the Father in Christ's name,
adore the one God in the distinction of Persons and equality of
majesty. The allegorical movement gathers Christ's people; the moral
movement teaches obedient trust and responsible praise; the anagogical
movement opens toward resurrection, secure fellowship and vision.
The child returns to health among his family, and the exiles
still long for home. Their different conditions can belong to one
hope because the giver is greater than the good already received.
